\documentclass[11pt,a4paper]{article}
\usepackage[margin=25mm]{geometry}
\usepackage{amsmath,amssymb,amsthm,hyperref}
\setlength{\emergencystretch}{2em}
\newtheorem{lemma}{Lemma}
\newtheorem{theorem}[lemma]{Theorem}
\newcommand{\gen}[1]{\langle#1\rangle}
\newcommand{\Sym}{\operatorname{Sym}}
\newcommand{\Alt}{\operatorname{Alt}}
\title{Even anchored factors and the remaining critical overlap}
\author{Human Compression project: a component of the L019 recovery}
\date{}
\begin{document}\maketitle
An endpoint anchor is a point moved by the first mark and fixed by
all later marks, or the reversed version of this condition. A marking
is fractured when every maximal omission is intransitive. All actions,
supports and intersections below are actual permutation actions.

\begin{lemma}\label{connected}
A transitive string marking with an endpoint anchor has connected
noncommutation diagram.
\end{lemma}
\begin{proof}
Reverse the order so the anchor $a$ belongs to the first mark.
Let $H$ be the group generated by the connected component containing
that mark. All other components commute with $H$ and fix $a$.
Consequently the full orbit of $a$ equals $a^H$, so $H$ is transitive.
An element commuting with a transitive group and fixing one point
fixes every point. Every mark outside that component would therefore
be the identity, contrary to the definition of a marking.
\end{proof}

\begin{theorem}\label{even}
Let a transitive fractured string C-marking of degree $d\ge12$
have an endpoint anchor, and suppose every mark is even.
Its rank $s$ satisfies $2s<d$.
\end{theorem}
\begin{proof}
We acquire two exact conclusions from the complete L044 recovery.
First, its unique-crossing theorem is stated for a transitive
connected string C-marking $G\le\Alt_d$ with all maximal omissions
intransitive. If some colour has exactly one eligible edge, it gives
$s\le\lfloor(d-1)/2\rfloor$. This statement does not require
$G=\Alt_d$. Its explicit smaller-degree alternating rank premises
are now supplied by the completed alternating rank theorem, and its
retained finite and attachment inputs remain part of the acquisition.
Lemma~\ref{connected} supplies its connectedness hypothesis here.

Suppose instead that no colour has exactly one eligible edge. Every
colour has at least one: otherwise its mark preserves every orbit
of its intransitive maximal omission, and adjoining it could not
produce a transitive group. Thus every colour has two eligible edges.
The acquired paired-cut count gives $2s\le d$. If equality holds,
the acquired square normal form and equality-connectedness proof
supply a connected selection with a single alternating square.
The complete ladder proof applies since $d>8$.

We use its full generator description, rather than only its invariant
pair partition. After possibly reversing the labels, the pairs are
\[
                    L,R_0,R_1,\ldots,R_{s-2}.
\]
The mark $\rho_0$ interchanges the two points of every $R_j$ and
may also interchange the two points of $L$. For $1\le i\le s-1$,
$\rho_i$ swaps the two adjacent pairs in the displayed path, by
two edges, and fixes all other points. The first such adjacent pair
is $L,R_0$. These are the full actual actions, including all edges
not originally selected.

Every point of an $R_j$ is consequently moved by $\rho_0$ and by
at least one mark of positive index. A point of $L$ is moved by
$\rho_1$, and possibly by $\rho_0$. If it is moved by just one
mark, that mark has index $1$. Since $s=d/2\ge6$, this is neither
endpoint index. The same conclusion holds after reversing labels.
There is no endpoint anchor, a contradiction. Equality is excluded;
both branches give the required strict inequality.
\end{proof}

For the next statement call an internal perfect split \emph{sufficient}
if its mark is pure, or if one internal factor is trivial and both
attaching groups contain an odd permutation. The acquired pure-leaf
and odd-attachment theorems prove that such a split admits the required
subdivision. This definition does not assert that every usable split
is sufficient.

\begin{lemma}[Transport through a one-point intersection]\label{transport}
Suppose the labels of a transitive marking are the concatenation of
two nonempty submarkings $P,E$, supported and transitive on $X,Y$,
where $X\cap Y=\{a\}$. Suppose all marks of $P$ except its last
fix $a$, and all marks of $E$ except its first fix $a$.
Every internal sufficient perfect split of either factor is an
internal sufficient perfect split of the whole marking. If the
whole marking is fractured, both factors are fractured.
\end{lemma}
\begin{proof}
Consider a perfect split inside $P$, with omission orbits $C,D$.
Its suffix contains the last mark, which must move $a$ by
transitivity of $P$. Hence $a\in D$: every prefix mark fixes $a$,
while the suffix is supported on $D$. The actual prefix remains
supported and transitive on $C$ in the whole marking. The old
suffix and $E$ are transitive on $D$ and $Y$ respectively, with
common point $a$, so their union is transitive on $D\cup Y$.
They fix $C$. The split mark fixes $Y\setminus\{a\}$ and has
exactly its old crossing and internal factors. The split is therefore
perfect, and purity and a trivial internal factor are preserved.
Each of the new attaching groups contains its old counterpart;
oddness is preserved too. Reverse the full order for a split in $E$.

If a maximal omission inside $P$ were transitive on $X$, adjoining
all of $E$ would make the corresponding whole omission transitive
on $X\cup Y$. This contradicts fracture. The argument for $E$
is identical. Faithfulness and the intersection property of each
factor are inherited on its actual support, without a quotient.
\end{proof}

\begin{theorem}[Conditional reduction to small even factors]\label{critical}
Let a fractured natural symmetric C-marking have degree $n\ge5$
and rank $r$ with $2r>n+2$. Assume that every such marking of
strictly smaller degree has an internal sufficient split.
Then the given marking has an internal sufficient split, or, after
reversal, its marks are the concatenation of two actual factors
on sets meeting in one point: a natural symmetric factor of even
degree $m$ and rank $m/2+1$, and an even transitive factor of degree
\[
                         e\in\{6,8,10\}
\]
and rank $e/2$. Both factors are fractured and have an anchor at
their common point. In the latter alternative $n=m+e-1$ and
$2r=n+3$.
\end{theorem}
\begin{proof}
The acquired trivial-factor theorem supplies an internal perfect
split. Reverse it to write
\[
 q=(a\ b)\beta,\quad L=G_{<i},\quad R=G_{>i},\quad
 A\sqcup B=\Omega,
\]
where $L,R$ are supported and transitive on $A,B$, and $\beta$
is supported in $B\setminus\{b\}$. Put $J=\gen{q,R}$.
It is supported and transitive on $\{a\}\cup B$, meeting the
support of $L$ exactly at $a$. The string relations and transitivity
give the endpoint anchors required by Lemma~\ref{transport}.
At least one of $L,J$ contains an odd permutation since together
they generate the natural symmetric group. If both do, the split
is sufficient. A pure $q$ also finishes immediately.

If $J$ is even, the acquired even-attachment theorem says that
$L=\Sym(A)$ and either $2i>|A|+2$, or $|A|=m$ is even,
$i=m/2+1$, and $J$ is the degree-six rank-three copy of $A_5$.
In the strict case induction inside the smaller factor $L$ and
Lemma~\ref{transport} finish. The exception has exactly the stated
form with $e=6$.

It remains that $L$ is even and $J$ is odd. The actual-wing theorem
makes $R=\Sym(B)$. Thus $\beta\in R$, so $q\beta=(a\ b)$
and $J=\Sym(\{a\}\cup B)$. Set $e=|A|$, $m=|B|+1$,
$i=\operatorname{rank}(L)$, and $k=\operatorname{rank}(J)$.
The nonsymmetric-wing bound gives $2i\le e$. Since $r=i+k$,
\[
                    2k>m+1+(e-2i)\ge m+1.
\]
If $2k>m+2$, induction inside the smaller actual symmetric factor
$J$ and Lemma~\ref{transport} finish. Otherwise integrality forces
$m$ even and $2k=m+2$. The same strict inequality now gives
$2i>e-1$, hence $2i=e$ and $e$ is even.

The factor $L$ is fractured by Lemma~\ref{transport} and has an
endpoint anchor. Theorem~\ref{even} excludes $e\ge12$.
Degree two has no nonidentity even involution. In degree four every
nonidentity even involution is fixed-point-free; since $i=2$,
the second mark cannot fix an anchor of the first. This excludes
$e=2,4$, leaving $6,8,10$. The degree and rank formulas in the
statement follow by counting the common point and the disjoint
label intervals. Every invocation of the inductive hypothesis has
strictly smaller actual degree; no critical case has been assumed
solved in this reduction.
\end{proof}

\paragraph{Attribution and remaining obligation.}
The alternating rank and paired-fracture results remain attributed
to Cameron--Fernandes--Leemans--Mixer through the complete L044
recovery. The new use here is the endpoint-anchor consequence of
the full ladder description. The one-point transport and the
parity accounting reduce the extension problem in
Cameron--Fernandes--Leemans, arXiv:2212.12723v2, to the explicitly
retained small even factors. They do not settle those critical
attachments or the reversible counting correspondence. No new
census, whole-original credit, conjecture, Lean theorem or compression
measurement is claimed.
\end{document}
